\begin{recipe}{Pastinaaksoep met croutons}
\kicker[Voorgerecht,Soep]{Voorgerecht · soep}
\index[register]{Pastinaaksoep met croutons}
\index[register]{Pastinaak}
\index[register]{Aardappel}

\meta[Nederlands]{35 minuten \dvd Voor 4 personen \dvd Nederlands}

\lettrine{D}{eze} romige pastinaaksoep maakten we voor kerst 2025. Niet moeilijk
om te maken, en met de zelfgebakken croutons en lente-ui erbovenop ziet hij er
op tafel meteen feestelijk uit.

\blockrule{Ingrediënten}
\begin{ingredients}
  \ing{500 g}{pastinaken}\index[register]{Pastinaak}
  \ing{500 g}{vastkokende aardappels}\index[register]{Aardappel}
  \ing{1}{knoflookteentje}
  \ing{1 bosje}{lente-ui}
  \ing{2 el}{olie}
  \ing{1 l}{groentebouillon}
  \ing{150 g}{crème fraîche}
  \ingb{snufje zout}
\end{ingredients}

\noindent\textit{Croutons:}
\begin{ingredients}
  \ing{4 plakken}{geroosterd brood}
  \ing{2 el}{boter}
  \ingb{beetje zout}
\end{ingredients}

\blockrule{Bereiding}
\begin{steps}
  \item Was de pastinaken en aardappelen, schil ze en snijd ze in blokjes.
  \item Pel de knoflooktenen en snijd ze fijn.
  \item Was de lente-uitjes, verwijder de worteluiteinden en snijd ze in fijne
        ringen.
  \item Verwarm de olie in een pan op middelhoog vuur en bak de pastinaken,
        aardappelen, knoflook en twee derde van de lente-uitjes ongeveer
        3 minuten.
  \item Blus af met de groentebouillon en laat het geheel ongeveer 20 minuten
        onder deksel garen.
  \item Snijd het geroosterde brood grof in blokjes.
  \item Smelt de boter in een koekenpan en rooster de broodblokjes ongeveer
        4 minuten tot ze goudbruin zijn.
  \item Laat de broodblokjes uitlekken op keukenpapier en bestrooi ze licht
        met zout.
  \item Voeg de crème fraîche toe aan de soep en pureer alles fijn.
  \item Breng op smaak, verdeel de soep over de borden en garneer met
        croutons en de resterende lente-ui.
\end{steps}
\end{recipe}
